\chapter*{Résumé}

La gestion de la maintenance des équipements aéroportuaires constitue un enjeu stratégique pour garantir la sécurité, la disponibilité et la performance des infrastructures. Face aux limites des approches correctives traditionnelles et à l'absence d'un outil intégré de suivi intelligent, ce projet de stage propose la conception et la réalisation de la plateforme \textbf{AIMOS} (Airport Intelligent Maintenance Operations System).

AIMOS est une plateforme web de maintenance intelligente destinée aux aéroports marocains, développée dans le cadre d'un stage au sein de l'\textbf{ONDA} (Office National Des Aéroports). Elle couvre l'ensemble du cycle de maintenance : gestion des équipements et des capteurs, planification et suivi des interventions avec gammes de maintenance (checklists), surveillance en temps réel via un système d'alertes automatiques, et un module d'intelligence artificielle pour la maintenance prédictive.

L'architecture technique repose sur un backend Django REST Framework, un frontend React avec Vite, une base de données PostgreSQL, le tout orchestré via Docker Compose et déployé en production avec Nginx et Gunicorn. Quatre rôles utilisateurs ont été identifiés et implémentés : Administrateur, Superviseur, Responsable maintenance et Technicien.

Les résultats obtenus démontrent une plateforme fonctionnelle offrant des tableaux de bord différenciés par rôle, un système d'alertes automatiques basé sur les seuils des capteurs, un module de prédiction de pannes utilisant la régression linéaire et la détection d'anomalies par z-score, ainsi qu'un système de notifications en temps réel.

\vspace{0.5cm}
\noindent\textbf{Mots-clés :} Maintenance prédictive, Aéroport, IoT, Intelligence artificielle, Capteurs, Alertes automatiques, Django, React, Docker, GMAO.
